\section{Future Work}

The most immediate priority is making the Knowledge Graph self-maintaining. A continuous integration pipeline that periodically re-scrapes clinical guideline sources, re-indexes the vector database, and flags divergent rules for human review would allow the graph to track evolving evidence without manual reconstruction. Paired with a lightweight approval interface where new rule batches require clinician sign-off before entering the live graph, this would close the human-in-the-loop gap in the current automated pipeline.

With continuous maintenance in place, the natural next step is transitioning from offline retrospective label correction to live prospective decision support at triage. The KG pipeline could infer needed diagnostics from presenting symptoms before physician assessment, directly serving the MLMD's parallelization objective. Computational cost becomes the binding constraint: the GraphRAG pipeline's 175\,ms per-encounter latency is acceptable at current SickKids volumes but would require model quantization, vector caching, or distilled biomedical encoders to scale across multi-facility deployments.

At scale, the efficiency gains compound. A 10-percentage-point recall improvement across SickKids' approximately 80,000 annual ED visits would redirect roughly 8,000 test authorizations per year from the serialized to the parallelized pathway. The broader contribution is methodological. Asymmetric label noise from unobserved external clinical events is a systematic problem across fragmented hospital networks, not a SickKids-specific artifact. The Knowledge Graph and meta-learning framework developed here offers a deployable approach for addressing it without cross-institutional EHR integration, a dependency that current interoperability infrastructure cannot satisfy.
